\section{Finite-overlap gap criteria}%
\label{sec:peps_region_parent_gap_criterion}

We apply the projector expansion of
\cite[Section~IV.C.3]{Cirac2021Matrix}, local source lines 2170--2179,
to the canonical regional parent interactions on a finite graph.
The anticommutator estimate below is an explicit hypothesis.
Deriving it from properties of a boundary state is a separate question.

\begin{theorem}[Disjoint regional interactions commute]%
    \label{thm:peps_disjoint_region_interactions_commute}
    \lean{TNLean.PEPS.regionReplaceConfig,
        TNLean.PEPS.regionLocalTerm_mulVec_apply,
        TNLean.PEPS.regionReplaceConfig_apply_of_disjoint,
        TNLean.PEPS.regionReplaceConfig_comm_of_disjoint,
        TNLean.PEPS.regionLocalTerm_commute_of_disjoint,
        TNLean.PEPS.regionLocalTermES,
        TNLean.PEPS.regionLocalTermES_commute_of_disjoint,
        TNLean.PEPS.regionParentHamiltonianES,
        TNLean.PEPS.regionParentHamiltonianES_eq_sum}
    \leanok
    \uses{def:peps_region_parent_hamiltonian}
    Let $R,S$ be disjoint vertex regions and let $h_R,h_S$
    be arbitrary linear operators on their respective physical spaces.
    Their extensions to the full physical space commute.
\end{theorem}

\begin{proof}\leanok
    The extension of $h_R$ sums only the physical indices in $R$
    and retains all complementary labels. Replacements on disjoint
    regions commute. Interchanging the two finite sums therefore
    identifies the two compositions.
\end{proof}

\begin{theorem}[Gap from overlapping-pair estimates]%
    \label{thm:peps_finite_overlap_gap}
    \lean{TNLean.PEPS.regionParentHamiltonianES_norm_gap_of_quadraticForm,
        TNLean.PEPS.regionParentHamiltonianES_quadraticForm_of_overlap_anticommutator,
        TNLean.PEPS.regionParentHamiltonianES_norm_gap_of_overlap_anticommutator,
        TNLean.PEPS.regionParentHamiltonianES_norm_gap_of_overlap_of_comparison}
    \leanok
    \uses{thm:peps_region_projector_unit_gap,
        thm:peps_disjoint_region_interactions_commute}
    Let $P_i$ be the extended canonical parent projector for
    a region $R_i$, and let $H=\sum_i P_i$.
    Suppose that each region overlaps at most $m$ other regions,
    where $m\ge1$, and that $\varepsilon\ge0$ satisfies
    \[
        P_iP_j+P_jP_i\ge-\varepsilon(P_i+P_j)
    \]
    for every distinct overlapping pair. Then
    \[
        H^2\ge(1-m\varepsilon)H.
    \]
    If $m\varepsilon<1$, then every $\xi\perp\ker H$ satisfies
    \[
        \|H\xi\|\ge(1-m\varepsilon)\|\xi\|.
    \]
    If another positive parent with the same regional kernels satisfies
    $H'\ge\kappa H$ for some $\kappa>0$, its norm gap is at least
    $\kappa(1-m\varepsilon)$.
\end{theorem}

\begin{proof}\leanok
    Expand $H^2$. The diagonal terms sum to $H$ since $P_i^2=P_i$.
    A disjoint pair consists of commuting orthogonal projections,
    so its anticommutator is positive. Sum the assumed lower bounds
    over the remaining unordered pairs. Each projector occurs
    at most $m$ times, giving the quadratic-form inequality.
    The spectral theorem for a positive finite-dimensional operator
    then gives the asserted norm bound on the complement of its kernel.
    The comparison assertion follows from the common kernel and
    the corresponding lower bound between positive operators.
\end{proof}

\begin{corollary}[A gap independent of the volume]%
    \label{cor:peps_uniform_finite_overlap_gap}
    \lean{TNLean.PEPS.regionParentHamiltonianES_uniform_norm_gap_of_overlap_anticommutator}
    \leanok
    \uses{thm:peps_finite_overlap_gap}
    Consider an arbitrary family of finite graphs, regional canonical
    PEPS parents, and physical and bond dimensions.
    If the overlap bound $m$ and anticommutator bound $\varepsilon$
    hold throughout the family and satisfy $m\varepsilon<1$,
    then every parent in the family has norm gap at least the same
    positive number $1-m\varepsilon$.
\end{corollary}

\begin{proof}\leanok
    Apply the preceding estimate to each member of the family.
    The resulting constant does not depend on the graph or its volume.
\end{proof}

\begin{theorem}[Unit gap for commuting canonical parents]%
    \label{thm:peps_commuting_region_parent_unit_gap}
    \lean{TNLean.PEPS.regionParentHamiltonianES_unit_norm_gap_of_commute,
        TNLean.PEPS.regionParentHamiltonianES_unit_norm_gap_of_pairwise_disjoint}
    \leanok
    \uses{thm:peps_region_projector_unit_gap,
        thm:peps_disjoint_region_interactions_commute}
    If the extended canonical parent projectors commute pairwise,
    their sum satisfies $H^2\ge H$ and has norm gap at least one
    on $(\ker H)^\perp$. In particular, this holds for
    pairwise disjoint regions.
\end{theorem}

\begin{proof}\leanok
    Every cross term in the projector expansion is nonnegative.
    Thus $H^2\ge H$, and the preceding spectral argument applies.
\end{proof}
